\section{证明}
\label{app:proofs}

本附录给出正文各命题的完整证明。负二项分布 $\NB(r,p)$ 取“失败次数”参数化，概率母函数为
\begin{equation}
  G(s)=\frac{p^{r}}{[1-(1-p)s]^{r}},\qquad \E X=\frac{r(1-p)}{p},\quad \Var X=\frac{r(1-p)}{p^2}.
  \label{eq:pgf}
\end{equation}
取 $r=k$、$p=k/(k+R)$，由式~\eqref{eq:pgf} 即得式~\eqref{eq:offspring} 中的均值与方差。

\subsection{负二项分布的求和与稀疏}
\label{app:lemma}

引理：若 $X_1,\dots,X_n$ 独立同分布于 $\NB(k,p)$，则 $I=\sum_iX_i\sim\NB(nk,p)$；若再有 $C\mid I\sim\Bin(I,\rho)$，则 $C\sim\NB(nk,p')$，其中 $p'=p/[1-(1-p)(1-\rho)]$。

证明：独立和的概率母函数等于各项之积，故 $G_I(s)=G(s)^n$，即 $\NB(nk,p)$。二项稀疏对应复合 $G_C(s)=G_I(1-\rho+\rho s)$。记 $q=1-(1-p)(1-\rho)$，则
\begin{equation*}
  G_C(s)=\frac{p^{nk}}{[q-(1-p)\rho s]^{nk}}=\frac{(p/q)^{nk}}{\{1-[(1-p)\rho/q]\,s\}^{nk}}.
\end{equation*}
由于 $1-p/q=(q-p)/q=(1-p)\rho/q$，$G_C$ 正是 $\NB(nk,p/q)$ 的母函数。代入 $p=k/(k+R_t)$，得 $q=(k+\rho R_t)/(k+R_t)$，$p'=k/(k+\rho R_t)$，这就是式~\eqref{eq:report-dist}。式~\eqref{eq:relvar} 由式~\eqref{eq:report-dist} 的均值与方差直接相除得到。\qed

\subsection{\texorpdfstring{命题~\ref{prop:current} 的证明}{命题 1 的证明}}

记 $c$ 为 $C_t$ 的观测值，$\theta=\rho R_t$。由式~\eqref{eq:report-dist}，对数似然为
\begin{equation*}
  \ell(R_t)=\text{常数}+nk\log\frac{k}{k+\theta}+c\log\frac{\theta}{k+\theta}.
\end{equation*}
对 $R_t$ 求导并整理，得分函数为
\begin{equation}
  S(R_t)=\frac{c}{R_t}-\frac{(c+nk)\rho}{k+\rho R_t}
  =\frac{k\rho n}{R_t(k+\rho R_t)}(\widehat R_t-R_t).
  \label{eq:score}
\end{equation}

(i) 当 $c>0$ 时，由式~\eqref{eq:score}，$S$ 在 $R_t<\widehat R_t$ 时为正，在 $R_t>\widehat R_t$ 时为负，似然在 $\widehat R_t$ 处取唯一最大值。当 $c=0$ 时，$S(R_t)=-nk\rho/(k+\rho R_t)<0$，似然在 $(0,\infty)$ 上严格递减，最大值在边界 $R_t=0$，与 $\widehat R_t=0$ 一致。

(ii) 由式~\eqref{eq:report-dist}，$\E\widehat R_t=\rho nR_t/(\rho n)=R_t$，且
\begin{equation*}
  \Var(\widehat R_t)=\frac{\rho nR_t+\rho^2nR_t^2/k}{(\rho n)^2}=\frac{R_t}{\rho n}+\frac{R_t^2}{kn}=\frac{R_t(k+\rho R_t)}{k\rho n}.
\end{equation*}

(iii) Fisher 信息为得分函数的方差。由式~\eqref{eq:score} 与 (ii) 得
\begin{equation*}
  \mathcal I(R_t)=\frac{(k\rho n)^2}{[R_t(k+\rho R_t)]^2}\Var(\widehat R_t)=\frac{k\rho n}{R_t(k+\rho R_t)}=\frac{n}{R_t/\rho+{R_t}^2/k}.
\end{equation*}
因此 $\Var(\widehat R_t)=1/\mathcal I(R_t)$。式~\eqref{eq:score} 表明得分函数是 $\widehat R_t-R_t$ 的倍数，这正是 Cramér--Rao 不等式取等号的充要条件，对每个有限 $n$ 都成立。

(iv) 式~\eqref{eq:report-dist} 的概率函数可写为 $h(c)\exp[\eta c-\psi(\eta)]$，其中自然参数 $\eta=\log[\rho R_t/(k+\rho R_t)]$。$k$ 已知时这是单参数指数族；$R_t$ 取遍 $(0,\infty)$ 时，$\eta$ 取遍 $(-\infty,0)$，含开区间，所以 $C_t$ 是完备充分统计量。$\widehat R_t$ 是 $C_t$ 的函数且无偏，由 Lehmann--Scheffé 定理，它是一致最小方差无偏估计\citep{lehmann1998point}。\qed

\subsection{\texorpdfstring{命题~\ref{prop:average} 的证明}{命题 2 的证明}}

由 H4，给定 $n$ 时 $R_t$ 的均值与方差为 $\Rbar$ 与 $v_R$。由命题~\ref{prop:current}(ii)，$\E(\widehat R_t\mid n,R_t)=R_t$，所以 $\E(\widehat R_t\mid n)=\Rbar$。由全方差公式
\begin{equation*}
  \Var(\widehat R_t\mid n)=\E\{R_t/(\rho n)+R_t^2/(kn)\mid n\}+\Var(R_t\mid n)
  =\frac{\Rbar}{\rho n}+\frac{\Rbar^2+v_R}{kn}+v_R.
\end{equation*}
以 $m=\rho n$ 代入，前两项之和为 $A/m$，得式~\eqref{eq:var-average}。

(i) $A/m\to0$。(ii) 规模项份额为 $(A/m)/(v_R+A/m)=1/(1+mv_R/A)=1/(1+m/\mx)$。(iii) 翻倍降幅为 $[V(m)-V(2m)]/V(m)=[A/(2m)]/(v_R+A/m)=1/[2(1+m/\mx)]$。(iv) $V'(m)=-A/m^2<0$，$V''(m)=2A/m^3>0$；弹性为 $mV'(m)/V(m)=-(A/m)/(v_R+A/m)=-\mx/(m+\mx)$，在 $m=\mx$ 处为 $-1/2$。\qed

式~\eqref{eq:pooling} 的推导如下。$J$ 个环境项方差同为 $v_R$、两两相关系数平均为 $\bar c$ 时，其等权平均的方差为 $J^{-2}[Jv_R+J(J-1)\bar cv_R]=v_R[\bar c+(1-\bar c)/J]$。

\subsection{\texorpdfstring{命题~\ref{prop:reporting} 的证明}{命题 3 的证明}}

把式~\eqref{eq:var-average} 写成 $n$ 与 $\rho$ 的函数
\begin{equation*}
  V(n,\rho)=v_R+\frac{\Rbar}{\rho n}+\frac{\Rbar^2+v_R}{kn}.
\end{equation*}
(i) 令 $\rho\to1$，第二项趋于 $\Rbar/n>0$，第三项与 $\rho$ 无关且为正，故极限严格大于 $v_R$。(ii) 令 $n\to\infty$，后两项都趋于零。(iii) 由 $1/\rho=1+(1-\rho)/\rho$ 直接得到。(iv) 两个方向的边际效应为
\begin{equation*}
  -\frac{\partial V}{\partial\log n}=\frac{\Rbar}{\rho n}+\frac{\Rbar^2+v_R}{kn}=\frac{A}{m},\qquad
  -\frac{\partial V}{\partial\log\rho}=\frac{\Rbar}{\rho n}=\frac{\Rbar}{m}.
\end{equation*}
二者之比为 $A/\Rbar=1+\rho(\Rbar^2+v_R)/(k\Rbar)$。\qed

\subsection{\texorpdfstring{命题~\ref{prop:allocation} 的证明}{命题 4 的证明}}

设各期上一代人数均为 $n$，给定 $(R_1,\dots,R_T)$ 时各期报告数相互独立。记 $\bar R_T=T^{-1}\sum_tR_t$。由全方差公式
\begin{equation*}
  \frac{1}{T^2}\Var({\textstyle\sum_t}\widehat R_t)=\Var(\bar R_T)+\frac1{T^2}\sum_t\E\{R_t/(\rho n)+R_t^2/(kn)\}
  =\Var(\bar R_T)+\frac{A}{Tm}.
\end{equation*}
第二项等于 $A/B$，只依赖 $R_t$ 的边际分布，与相关结构无关。$R_t$ 独立时 $\Var(\bar R_T)=v_R/T$，得式~\eqref{eq:alloc-indep}。

一阶自回归情形下，$\Cov(R_s,R_t)=v_Rr^{|s-t|}$，于是
\begin{equation*}
  T^2\Var(\bar R_T)=v_R\,T+2v_R\sum_{h=1}^{T-1}(T-h)r^h,\qquad
  \sum_{h=1}^{T-1}(T-h)r^h=\frac{r[T(1-r)-(1-r^T)]}{(1-r)^2}.
\end{equation*}
后一个等式可对 $T$ 归纳验证。代入整理即得式~\eqref{eq:alloc-ar}。

单调性。记 $S_T=T^2\Var(\bar R_T)$，$\Delta_j=v_R(1+2\sum_{h=1}^{j}r^h)$，则 $S_{T+1}=S_T+\Delta_T$，$S_T=\sum_{j=0}^{T-1}\Delta_j$。$\Var(\bar R_{T+1})<\Var(\bar R_T)$ 等价于 $T^2\Delta_T<(2T+1)S_T$。数列 $\Delta_j$ 的增量 $2v_Rr^{j+1}$ 递减，所以它是凹数列，满足 $\Delta_j\geq[(T-j)\Delta_0+j\Delta_T]/T$。对 $j=0,\dots,T-1$ 求和，得 $S_T\geq[(T+1)\Delta_0+(T-1)\Delta_T]/2$。记 $d=\Delta_T-\Delta_0$，代入 $\Delta_0=v_R$，有
\begin{equation*}
  (2T+1)S_T-T^2\Delta_T\geq(T+1)(Tv_R-d/2).
\end{equation*}
由于 $d=2v_R\sum_{h=1}^{T}r^h<2Tv_R$，右端为正。$r=0$ 时即独立情形，单调性显然。\qed

\subsection{\texorpdfstring{推论~\ref{cor:optimal} 的证明}{推论 1 的证明}}

由式~\eqref{eq:alloc-indep}，以 $T=\mathcal B/(c_0+c_1m)$、$B=Tm$ 代入，得
\begin{equation*}
  f(m)=\frac{(c_0+c_1m)(v_R+A/m)}{\mathcal B}=\frac{c_0v_R+c_1A+c_1v_Rm+c_0A/m}{\mathcal B}.
\end{equation*}
$f'(m)=(c_1v_R-c_0A/m^2)/\mathcal B$，$f''(m)=2c_0A/(m^3\mathcal B)>0$。$c_0>0$ 时 $f$ 在 $(0,\infty)$ 上严格凸，唯一驻点为 $m^{*2}=c_0A/(c_1v_R)=(c_0/c_1)\mx$。代回得 $f(m^*)=[c_0v_R+c_1A+2\sqrt{c_0c_1v_RA}]/\mathcal B$，即式~\eqref{eq:optimal-m}。

约束 $T\geq1$ 等价于 $m\leq(\mathcal B-c_0)/c_1$。若 $m^*$ 超出这一范围，$f$ 在可行区间上递减，最优解取右端点，即 $T=1$。若 $c_0=0$，$f(m)=c_1(v_Rm+A)/\mathcal B$ 严格递增，下确界在 $m\to0$ 处取得。

整数解。$m=(\mathcal B/T-c_0)/c_1$ 关于 $T$ 单调，而 $f$ 关于 $m$ 单峰，所以 $f$ 作为 $T$ 的函数也单峰，整数最优解是连续最优期数两侧的两个整数之一。正文的取整例子即按此比较。环境服从一阶自回归且 $T$ 较大时，环境项约为 $v_R(1+r)/[(1-r)T]$，相当于把 $v_R$ 换成 $v_R(1+r)/(1-r)$，于是 $m^*$ 约乘以 $\sqrt{(1-r)/(1+r)}$。\qed

\subsection{\texorpdfstring{命题~\ref{prop:test} 的证明}{命题 5 的证明}}

取 $R>R'$。由式~\eqref{eq:report-dist}，两者的似然比为
\begin{equation*}
  \frac{f(c;R)}{f(c;R')}=\frac{(k+\rho R')^{nk}\,[R(k+\rho R')]^{c}}{(k+\rho R)^{nk}\,[R'(k+\rho R)]^{c}}.
\end{equation*}
由于 $R/(k+\rho R)$ 关于 $R$ 严格递增，方括号内大于 1，似然比关于 $c$ 严格递增，分布族具有单调似然比。

(i) 对单调似然比族，形如 $\{C_t\geq c\}$ 的检验的功效函数关于 $R_t$ 非降\citep{lehmann2005testing}，所以在 $R_t\leq1$ 上的最大拒绝概率在 $R_t=1$ 处取得，等于 $\alpha^*=P_{R_t=1}(C_t\geq c_\alpha)\leq\alpha$。

(ii) 由 Karlin--Rubin 定理\citep{karlin1956theory}，对任意水平，形如“$C_t>c$ 时拒绝，$C_t=c$ 时以概率 $\gamma$ 拒绝”的检验在该水平下一致最优。非随机检验 $\{C_t\geq c_\alpha\}$ 是 $c=c_\alpha-1$、$\gamma=1$ 的特例，所以在水平 $\alpha^*$ 下一致最优；适当选取 $\gamma$ 可以使水平恰为 $\alpha$。

(iii) $R_t=1+\delta$ 时，$C_t\sim\NB(nk,k/[k+\rho(1+\delta)])$，拒绝概率为 $1-F_{\NB}(c_\alpha-1)$，即式~\eqref{eq:power}。\qed

\subsection{\texorpdfstring{注~\ref{rem:information}、式~\eqref{eq:mse-decomp} 与式~\eqref{eq:b-delta} 的推导}{注 1、式 (12) 与式 (15) 的推导}}

以下所有期望都在 $\mathcal F_t$ 条件下取，$\Rbar$、$v_R$ 理解为 $R_{t+1}$ 的条件均值与方差。记 $\bar\Lambda=\E(\Lambda_{t+1}\mid\mathcal F_t)$，$m=\rho\bar\Lambda$。第一次使用全方差公式，以 $(\Lambda_{t+1},R_{t+1})$ 为条件
\begin{equation*}
  \Var(C_{t+1})=\E(\rho\Lambda_{t+1}R_{t+1}+\rho^2\Lambda_{t+1}R_{t+1}^2/k)+\Var(\rho\Lambda_{t+1}R_{t+1}).
\end{equation*}
由条件独立，第一项为 $\rho\bar\Lambda\Rbar+\rho^2\bar\Lambda(\Rbar^2+v_R)/k=m\,A$。第二步计算乘积的方差
\begin{equation*}
  \Var(\Lambda_{t+1}R_{t+1})=\E(\Lambda_{t+1}^2)\E(R_{t+1}^2)-\bar\Lambda^2\Rbar^2
  =\bar\Lambda^2[v_R+\CV^2(\Lambda_{t+1})(\Rbar^2+v_R)].
\end{equation*}
两项相加后除以 $m^2$，即得式~\eqref{eq:info-decomp}。

式~\eqref{eq:mse-decomp} 的推导如下。$D_t$ 由 $\mathcal F_t$ 决定，所以 $\E(G_t\mid\mathcal F_t)=\mu_t/D_t$，$\Var(G_t\mid\mathcal F_t)=(\sigma_t/D_t)^2$。对 $\E[(G_t-\Rbar)^2\mid\mathcal F_t]=\Var(G_t\mid\mathcal F_t)+[\E(G_t\mid\mathcal F_t)-\Rbar]^2$ 再取期望即可。式~\eqref{eq:identity} 是两边同除以 $\Rbar D_t$ 的代数恒等式。

式~\eqref{eq:b-delta} 的推导换用另一个条件。记 $\mathcal H_t$ 为截至第 $t$ 周的真实感染史，并把式~\eqref{eq:dt} 的滞后权重当作感染压力的权重，即 $\Lambda_{t+1}=\sum_s w_sI_{t+1-s}$，$\sum_s w_s=1$。给定 $\mathcal H_t$，$\Lambda_{t+1}$ 已经确定，取 $m=\rho\Lambda_{t+1}$。它与正文的 $m=\rho\E(\Lambda_{t+1}\mid\mathcal F_t)$ 相差一层报告噪声，式~\eqref{eq:affine} 的拟合又以 $D_t$ 代替 $m$，式~\eqref{eq:b-delta} 描述的正是这层噪声。各周报告独立地服从 $C_{t+1-s}\mid I_{t+1-s}\sim\Bin(I_{t+1-s},\rho)$，所以分母的条件均值与条件方差为
\begin{equation*}
  \E(D_t\mid\mathcal H_t)=\rho\Lambda_{t+1},\qquad
  \Var(D_t\mid\mathcal H_t)=\rho(1-\rho)\sum_s {w_s}^2I_{t+1-s}.
\end{equation*}
在平稳规模附近取 $I_{t+1-s}\approx\Lambda_{t+1}$，并令 $D_t=m(1+\delta)$，则 $\E\delta=0$，$\E\delta^2=c/m$，其中 $c=(1-\rho)\sum_s {w_s}^2$，而 $\E\delta^3$ 与 $\E\delta^4$ 都是 $O(m^{-2})$。记 $X=C_{t+1}/m$，于是 $G_t=X/(1+\delta)$。给定 $\mathcal H_t$，$C_{t+1}$ 的报告抽样与此前各周的报告抽样相互独立，$X$ 与 $\delta$ 独立，展开中没有交叉项。与第~\ref{sec:variance}~节相同，$\E X=\Rbar$，$\Var X=v_R+A/m$。把 $(1+\delta)^{-1}$ 展开到 $\delta^2$ 再取期望，得到
\begin{equation*}
  \E(G_t\mid\mathcal H_t)=\Rbar\Bigl(1+\frac cm\Bigr)+O(m^{-2}),\qquad
  \Var(G_t\mid\mathcal H_t)=v_R+\frac{A+(\Rbar^2+3v_R)\,c}{m}+O(m^{-2}).
\end{equation*}
只保留 $\delta$ 的一阶项时，方差的附加项是 $\E(R_{t+1}^2)\,c/m=(\Rbar^2+v_R)\,c/m$。$\delta^2$ 的期望与 $1/m$ 同阶，它另外贡献 $2v_Rc/m$，并使条件均值上移 $\Rbar c/m$。以条件均值的平方归一化，两处二阶效应相互抵消：
\begin{equation*}
  \frac{\Var(G_t\mid\mathcal H_t)}{[\E(G_t\mid\mathcal H_t)]^2}
  =\frac{v_R}{\Rbar^2}+\frac{A+(1-\rho)(\Rbar^2+v_R)\sum_s {w_s}^2}{\Rbar^2\,m}+O(m^{-2}).
\end{equation*}
上式中 $1/m$ 的系数就是式~\eqref{eq:b-delta}。若按 $L_t$ 的定义以固定的 $\Rbar$ 归一化，系数变为 $[A+(1-\rho)(\Rbar^2+3v_R)\sum_s {w_s}^2]/\Rbar^2$。在第~\ref{sec:empirical-inflation}~节的情景参数下，主规格因此多出 0.021，约为式~\eqref{eq:b-delta} 的 1.8\%，正文略去了这一项。

这一推导用到三个条件。第一，$\sum_s w_s=1$ 且 $I_{t+1-s}\approx\Lambda_{t+1}$。规模以周增长比 $r$ 变化时，$\sum_s {w_s}^2$ 应换成 $\sum_s {w_s}^2I_{t+1-s}/\Lambda_{t+1}$，主规格在 $r=1.67$ 时为 0.555，在 $r=0.8$ 时为 0.461。第二，各周报告抽样条件独立。若报告比例存在跨周的共同冲击，$X$ 与 $\delta$ 相关，展开中会多出交叉项。第三，$m$ 要足够大。对固定的 $m$ 逐点枚举二项分布并截断 $D_t\geq5$，主规格下 $m[\CV^2(G_t\mid\mathcal H_t)-v_R/\Rbar^2]$ 在 $m=20$ 时为 1.223，在 $m=320$ 时为 1.159，按 $1/m$ 外推为 1.155，与式~\eqref{eq:b-delta} 一致；只用当周权重时，$m=20$ 处为 1.972，高阶余项已不可忽略。实证拟合以 $D_t$ 代替 $m$，属于含噪回归元。第~\ref{sec:empirical-inflation}~节模拟值与预测值的差距来自这一步、$D_t\geq5$ 的截断和小规模处的高阶余项，三者方向不同，部分相互抵消。

\section{复现说明}
\label{app:repro}

第~\ref{sec:variance}~节的蒙特卡洛核对在 $m\in[1,1000]$ 的 10 个规模上各做 20,000 次模拟，随机种子 20261013。第~\ref{sec:renewal-sim}~节的更新过程仿真每个初始规模 400 条轨迹，随机种子 20261014。第~\ref{sec:design-test}~节的检验水平与功效全部由负二项分布函数精确计算，没有模拟误差。

实证数据为 NHSN 公开的州级每周流感住院数与医院填报比例\citep{cdc2026nhsn}，样本截至 2026 年 9 月。仿射拟合的区间由按州聚类的 1,500 次重抽得到，工具变量回归使用 10,000 次重抽。FluSight 与欧洲 COVID-19 数据分别取自文献\citep{mathis2024evaluation, flaxman2020estimating} 公开的数据，聚类单位分别为州和国家。全部仿真、拟合与作图脚本随中间文件一并保存。
